\chapter{INTRODUCTION}

\section{Background}
Air pollution is one of the major environmental and public-health concerns in urban areas. Fine particulate matter, commonly known as PM$_{2.5}$, is especially important because it consists of very small airborne particles that can enter deep into the respiratory system and affect human health. PM$_{2.5}$ concentration is influenced by many factors such as vehicular emissions, road dust, construction activities, industrial sources, seasonal variation, temperature, humidity, wind speed, wind direction, and local geography.

Kathmandu Valley is a suitable study area for PM$_{2.5}$ forecasting because it experiences frequent air pollution problems due to dense urban activity, traffic emissions, construction dust, and its bowl-shaped valley structure. The surrounding hills and limited air circulation can contribute to pollutant accumulation under certain meteorological conditions. Therefore, short-term PM$_{2.5}$ forecasting in Kathmandu Valley can support public awareness, environmental monitoring, and air quality interpretation.

Air quality forecasting is a spatio-temporal problem. This means that pollution changes across both location and time. The future PM$_{2.5}$ concentration at one monitoring station may depend not only on its own previous readings, but also on nearby stations, wind movement, and meteorological conditions. Traditional time-series models can learn temporal patterns, but they may not fully capture spatial influence between monitoring stations.

\glspl{gnn} provide a suitable approach for this type of problem because air quality monitoring stations can be represented as a graph. In this graph, each monitoring station is treated as a node, and the relationship between stations is represented as an edge. By combining graph-based spatial learning with recurrent temporal learning, the system can model both station-to-station influence and time-based pollution variation.

In this project, an explainable wind-aware spatio-temporal graph neural network framework is proposed for PM$_{2.5}$ forecasting in Kathmandu Valley. The system will use historical PM$_{2.5}$ observations, meteorological variables, and station relationships to forecast PM$_{2.5}$ concentration after 6 hours. It will also provide explanation outputs to show which stations, environmental features, and previous time patterns influenced the prediction.

\section{Motivation}
Existing air quality forecasting approaches often focus mainly on individual monitoring stations or use fixed spatial relationships between stations. However, pollution does not remain fixed in one location. Wind can transport polluted air from one area to another, and the influence between stations can change depending on wind direction and wind speed. A fixed graph based only on distance may therefore fail to represent the changing nature of pollutant movement.

Kathmandu Valley also has limited monitoring coverage compared to the complexity of its urban environment. Some areas may not have physical air quality sensors, even though people living in those areas still need information about local air quality. This motivates the use of graph-based modeling and virtual nodes, where selected unmonitored map locations can be estimated using nearby real monitoring stations and wind-aware spatial relationships.

Another motivation is explainability. A forecasting system should not only produce a predicted PM$_{2.5}$ value; it should also provide understandable reasons behind the prediction. For example, the system should help explain whether a forecast was influenced by high PM$_{2.5}$ at a nearby station, wind flow from a polluted direction, low wind speed, humidity, or recent pollution trends. Such explanations can make the system more transparent and useful for environmental monitoring, interpretation, and public awareness.

\section{Problem Definition}
Short-term PM$_{2.5}$ forecasting in Kathmandu Valley is challenging because pollutant concentration is affected by both temporal and spatial factors. PM$_{2.5}$ at a particular station may depend on its previous values, meteorological conditions, wind direction, wind speed, and pollutant levels at nearby or upwind stations. Many conventional forecasting methods treat stations independently or use static spatial relationships, which may not properly represent changing pollutant transport caused by wind.

In addition, air quality monitoring stations are limited in number and cannot cover every road, neighborhood, or residential area. This creates difficulty in estimating PM$_{2.5}$ concentration at unmonitored locations. Many forecasting systems also provide only numerical predictions without explaining why the model made a certain forecast.

Therefore, the problem addressed in this project is to design and develop an explainable wind-aware spatio-temporal graph neural network model that can forecast PM$_{2.5}$ concentration after 6 hours for selected Kathmandu Valley monitoring stations and selected virtual locations. The system will model station relationships using distance, PM$_{2.5}$ correlation, wind direction, and wind speed, and will present interpretable outputs through a dashboard.

\section{Objectives}
The objectives of this project are:
\begin{enumerate}
    \item To build a wind-aware dynamic graph from PM$_{2.5}$ and meteorological data, incorporating distance, correlation, wind factors, and virtual nodes for spatial modeling.
    \item To develop and evaluate GAT-GRU and baseline models for PM$_{2.5}$ forecasting with an explainable dashboard showing forecasts, AQI, spatial patterns, virtual nodes, and model insights.
\end{enumerate}


\section{Scope and Applications}
The scope of this project focuses on PM$_{2.5}$ forecasting in Kathmandu Valley using selected monitoring stations from Kathmandu, Lalitpur, and Bhaktapur. The system will use historical PM$_{2.5}$ observations and meteorological variables such as wind speed, wind direction, temperature, humidity, pressure, dew point, and temporal features where available.

The project will represent monitoring stations as graph nodes and model the relationship between stations using distance, PM$_{2.5}$ correlation, and wind-aware features. In addition to physical monitoring stations, selected unmonitored map locations may be represented as virtual nodes to estimate PM$_{2.5}$ levels in areas without direct sensor readings.

The system will be developed as a research prototype for forecasting, explanation, and visualization. The major applications include PM$_{2.5}$ forecasting, station-level air quality interpretation, virtual-node-based estimation for unmonitored locations, pollution risk visualization, explainable environmental monitoring, and academic research support. The dashboard will help users understand predicted PM$_{2.5}$ levels, \gls{aqi} categories, spatial pollution patterns, important influencing stations, and important environmental features behind each prediction.
